% GENERATED FILE. Edit canonical lessons, then run npm run generate:books.
% canonical-source-hash: fnv1a64:678b8accd02fe6e8
% canonical-lessons: TA-C330-peraracu, TA-C330-cakkaravartti, TA-C330-karuvulam, TA-C330-kappam, TA-C330-onku

\chapter{An Empire}
\label{ch:ta-an-empire}
\hlchaptermodality{330}

\begin{chapteropening}
\textbf{By the end of this chapter:} \emph{I can talk about empires and their wealth.}

Complete the last lesson of chapter 330: I can talk about empires and their wealth.
\end{chapteropening}

\section[\texorpdfstring{\ta{பேரரசு} (pēraracu)}{பேரரசு (pēraracu)}]{\ta{பேரரசு} (pēraracu) — an empire}

\label{lesson:TA-C330-peraracu}



\begin{quote}
Before the new one: say the Tamil for a reign, then the Tamil for an envoy.
\end{quote}

\subsection*{You'll want to know: \ta{பேரரசு}}
\textbf{\ta{பேரரசு}} — \emph{pēraracu} — \textquotedblleft{}an empire\textquotedblright{}.

\begin{quote}
\ta{அந்தப்} \ta{பேரரசு} \ta{கடல்} \ta{வரை} \ta{பரவியது}.
\end{quote}

\begin{quote}
\emph{antap pēraracu kaṭal varai paraviyatu.}
\end{quote}

\begin{quote}
That empire spread as far as the sea.
\end{quote}

\ta{பேர்}, great + \ta{அரசு}, kingdom.

\begin{cousinweb}
\textbf{\ta{பேரரசு}} (\emph{pēraracu}) is a Tamil compound built from \ta{பேர்}, \textquotedblleft{}great\textquotedblright{} (DEDR 4411), and \ta{அரசு}, \textquotedblleft{}kingdom\textquotedblright{}, which DEDR 201 traces to Sanskrit \emph{rājan}.
\end{cousinweb}

\begin{grammarlens}[title={What you've built}]
One more word for telling what happened.
\end{grammarlens}

\subsection*{Your turn}
\begin{itemize}
\raggedright
  \item \emph{Say it:} \emph{pēraracu}
  \item \emph{Say it:} \emph{pēraracu}, once more
  \item \emph{Say it:} \emph{antap pēraracu kaṭal varai paraviyatu.} — That empire spread as far as the sea.
  \item \emph{Recall:} say the Tamil for a reign, then the Tamil for an envoy, then say \emph{pēraracu} again
\end{itemize}

\begin{tcolorbox}[breakable,colback=teal!4,colframe=teal!35!black,title={Before you move on}]
What does \ta{பேரரசு} mean? (\textquotedblleft{}An empire\textquotedblright{}.) Say it once more.
\end{tcolorbox}
\section[\texorpdfstring{\ta{சக்கரவர்த்தி} (cakkaravartti)}{சக்கரவர்த்தி (cakkaravartti)}]{\ta{சக்கரவர்த்தி} (cakkaravartti) — an emperor}

\label{lesson:TA-C330-cakkaravartti}



\begin{quote}
Before the new one: say the Tamil for an envoy, then the Tamil for an empire.
\end{quote}

\subsection*{You'll want to know: \ta{சக்கரவர்த்தி}}
\textbf{\ta{சக்கரவர்த்தி}} — \emph{cakkaravartti} — \textquotedblleft{}an emperor\textquotedblright{}.

\begin{quote}
\ta{சக்கரவர்த்தி} \ta{பல} \ta{நாடுகளை} \ta{ஆண்டார்}.
\end{quote}

\begin{quote}
\emph{cakkaravartti pala nāṭukaḷai āṇṭār.}
\end{quote}

\begin{quote}
The emperor ruled many lands.
\end{quote}

The ruler of a \ta{பேரரசு}.

\begin{cousinweb}
\textbf{\ta{சக்கரவர்த்தி}} (\emph{cakkaravartti}) is borrowed from Sanskrit \emph{cakravartin}, \textquotedblleft{}a universal ruler\textquotedblright{} (CDIAL 4549).
\end{cousinweb}

\begin{grammarlens}[title={What you've built}]
One more word for telling what happened.
\end{grammarlens}

\subsection*{Your turn}
\begin{itemize}
\raggedright
  \item \emph{Say it:} \emph{cakkaravartti}
  \item \emph{Say it:} \emph{cakkaravartti}, once more
  \item \emph{Say it:} \emph{cakkaravartti pala nāṭukaḷai āṇṭār.} — The emperor ruled many lands.
  \item \emph{Recall:} say the Tamil for an envoy, then the Tamil for an empire, then say \emph{cakkaravartti} again
\end{itemize}

\begin{tcolorbox}[breakable,colback=teal!4,colframe=teal!35!black,title={Before you move on}]
What does \ta{சக்கரவர்த்தி} mean? (\textquotedblleft{}An emperor\textquotedblright{}.) Say it once more.
\end{tcolorbox}
\section[\texorpdfstring{\ta{கருவூலம்} (karuvūlam)}{கருவூலம் (karuvūlam)}]{\ta{கருவூலம்} (karuvūlam) — a treasury}

\label{lesson:TA-C330-karuvulam}



\begin{quote}
Before the new one: say the Tamil for an empire, then the Tamil for an emperor.
\end{quote}

\subsection*{You'll want to know: \ta{கருவூலம்}}
\textbf{\ta{கருவூலம்}} — \emph{karuvūlam} — \textquotedblleft{}a treasury\textquotedblright{}.

\begin{quote}
\ta{மன்னரின்} \ta{கருவூலத்தில்} \ta{நிறையத்} \ta{தங்கம்} \ta{இருந்தது}.
\end{quote}

\begin{quote}
\emph{maṉṉariṉ karuvūlattil niṟaiyat taṅkam iruntatu.}
\end{quote}

\begin{quote}
There was a lot of gold in the king's treasury.
\end{quote}

Today also a government treasury.

\begin{cousinweb}
\textbf{\ta{கருவூலம்}} (\emph{karuvūlam}), \textquotedblleft{}a treasury\textquotedblright{}, is a Tamil word whose origin is not settled in the sources at hand.
\end{cousinweb}

\begin{grammarlens}[title={What you've built}]
One more word for telling what happened.
\end{grammarlens}

\subsection*{Your turn}
\begin{itemize}
\raggedright
  \item \emph{Say it:} \emph{karuvūlam}
  \item \emph{Say it:} \emph{karuvūlam}, once more
  \item \emph{Say it:} \emph{maṉṉariṉ karuvūlattil niṟaiyat taṅkam iruntatu.} — There was a lot of gold in the king's treasury.
  \item \emph{Recall:} say the Tamil for an empire, then the Tamil for an emperor, then say \emph{karuvūlam} again
\end{itemize}

\begin{tcolorbox}[breakable,colback=teal!4,colframe=teal!35!black,title={Before you move on}]
What does \ta{கருவூலம்} mean? (\textquotedblleft{}A treasury\textquotedblright{}.) Say it once more.
\end{tcolorbox}
\section[\texorpdfstring{\ta{கப்பம்} (kappam)}{கப்பம் (kappam)}]{\ta{கப்பம்} (kappam) — a royal tribute}

\label{lesson:TA-C330-kappam}



\begin{quote}
Before the new one: say the Tamil for an emperor, then the Tamil for a treasury.
\end{quote}

\subsection*{You'll want to know: \ta{கப்பம்}}
\textbf{\ta{கப்பம்}} — \emph{kappam} — \textquotedblleft{}a royal tribute\textquotedblright{}.

\begin{quote}
\ta{சிறிய} \ta{நாடுகள்} \ta{சக்கரவர்த்திக்குக்} \ta{கப்பம்} \ta{கட்டின}.
\end{quote}

\begin{quote}
\emph{ciṟiya nāṭukaḷ cakkaravarttikkuk kappam kaṭṭiṉa.}
\end{quote}

\begin{quote}
The small kingdoms paid tribute to the emperor.
\end{quote}

\ta{கப்பம்} \ta{கட்டு}, to pay tribute.

\begin{cousinweb}
\textbf{\ta{கப்பம்}} (\emph{kappam}) has its own entry in DEDR (1218) across the Dravidian languages, which also compares an Apabhraṃśa \emph{kappa}, \textquotedblleft{}tribute\textquotedblright{}, so its ultimate origin is not settled.
\end{cousinweb}

\begin{grammarlens}[title={What you've built}]
One more word for telling what happened.
\end{grammarlens}

\subsection*{Your turn}
\begin{itemize}
\raggedright
  \item \emph{Say it:} \emph{kappam}
  \item \emph{Say it:} \emph{kappam}, once more
  \item \emph{Say it:} \emph{ciṟiya nāṭukaḷ cakkaravarttikkuk kappam kaṭṭiṉa.} — The small kingdoms paid tribute to the emperor.
  \item \emph{Recall:} say the Tamil for an emperor, then the Tamil for a treasury, then say \emph{kappam} again
\end{itemize}

\begin{tcolorbox}[breakable,colback=teal!4,colframe=teal!35!black,title={Before you move on}]
What does \ta{கப்பம்} mean? (\textquotedblleft{}A royal tribute\textquotedblright{}.) Say it once more.
\end{tcolorbox}
\section[\texorpdfstring{\ta{ஓங்கு} (ōṅku)}{ஓங்கு (ōṅku)}]{\ta{ஓங்கு} (ōṅku) — to flourish; to rise high}

\label{lesson:TA-C330-onku}



\begin{quote}
Before the new one: say the Tamil for a treasury, then the Tamil for a royal tribute.
\end{quote}

\subsection*{You'll want to know: \ta{ஓங்கு}}
\textbf{\ta{ஓங்கு}} — \emph{ōṅku} — \textquotedblleft{}to flourish; to rise high\textquotedblright{}.

\begin{quote}
\ta{அவருடைய} \ta{ஆட்சியில்} \ta{கலைகள்} \ta{ஓங்கின}.
\end{quote}

\begin{quote}
\emph{avaruṭaiya āṭciyil kalaikaḷ ōṅkiṉa.}
\end{quote}

\begin{quote}
The arts flourished under his reign.
\end{quote}

Past (it …): \ta{ஓங்கியது}, \emph{ōṅkiyatu}.

Literally to rise high: \ta{ஓங்கி} \ta{வளர்ந்த} \ta{மரம்}, a tall tree.

\begin{cousinweb}
\textbf{\ta{ஓங்கு}} (\emph{ōṅku}) is native Dravidian (DEDR 1033): to grow, rise high, be exalted.
\end{cousinweb}

\begin{grammarlens}[title={What you've built}]
One more word for telling what happened.
\end{grammarlens}

\subsection*{Your turn}
\begin{itemize}
\raggedright
  \item \emph{Say it:} \emph{ōṅku}
  \item \emph{Say it:} \emph{ōṅku}, once more
  \item \emph{Say it:} \emph{avaruṭaiya āṭciyil kalaikaḷ ōṅkiṉa.} — The arts flourished under his reign.
  \item \emph{Recall:} say the Tamil for a treasury, then the Tamil for a royal tribute, then say \emph{ōṅku} again
\end{itemize}

\begin{tcolorbox}[breakable,colback=teal!4,colframe=teal!35!black,title={Before you move on}]
What does \ta{ஓங்கு} mean? (\textquotedblleft{}To flourish; to rise high\textquotedblright{}.) Say it once more.
\end{tcolorbox}
